\documentclass[12pt,a4paper]{article}

\usepackage[utf8]{inputenc}
\usepackage[T1]{fontenc}
\usepackage{amsmath}
\usepackage{amssymb}
\usepackage{physics}
\usepackage{booktabs}
\usepackage{array}
\usepackage{geometry}
\usepackage{hyperref}
\usepackage{cite}

\geometry{
  a4paper,
  top=2.5cm,
  bottom=2.5cm,
  left=2.5cm,
  right=2.5cm
}

\hypersetup{
  colorlinks=true,
  linkcolor=blue,
  citecolor=blue,
  urlcolor=blue
}

\title{\textbf{Coulomb-Barrier-Free Neutron Production via\\
Proton--Electron Capture in Dense Hydrogen Gas\\
with Lithium-6 Energy Extraction}}

\author{Eymen Aydoğan}

\date{May 2026}

\begin{document}

\maketitle

\begin{abstract}
We propose a theoretical mechanism for neutron production and energy
extraction that entirely avoids the Coulomb barrier. A free proton
injected into dense hydrogen gas interacts with a bound electron via
electrostatic attraction, producing a neutron through inverse beta
decay ($p + e^- \rightarrow n + \nu_e$). The liberated hydrogen
nucleus then becomes a free proton, capable of repeating the reaction
on a neighboring hydrogen atom, forming a self-propagating chain.
Produced neutrons are captured by Lithium-6, releasing 4.78 MeV per
event via $n + {}^6\text{Li} \rightarrow {}^4\text{He} + {}^3\text{H}
+ 4.78\,\text{MeV}$. The net energy output per cycle is 3.487 MeV.
The mechanism requires no Coulomb barrier crossing and represents a
novel conceptual pathway for low-energy neutron-driven nuclear energy.
\end{abstract}

\noindent\textbf{Keywords:} inverse beta decay, Coulomb barrier,
chain reaction, neutron capture, lithium-6, LENR, hydrogen gas

\hrule
\vspace{1em}

\section{Introduction}

The Coulomb barrier is the primary obstacle in achieving controlled
nuclear reactions at low energies. In conventional fusion, two
positively charged nuclei must overcome a repulsive potential on the
order of 1 MeV at nuclear distances, requiring extreme temperatures
or pressures to proceed at useful rates.

We propose a fundamentally different approach: instead of fusing two
positive nuclei, we exploit the \emph{attractive} interaction between
a free proton and a bound electron in a hydrogen atom. Since the
proton and electron carry opposite charges, no repulsive Coulomb
barrier exists. The resulting neutrons, being electrically neutral,
can penetrate target nuclei without any Coulomb barrier, making
Lithium-6 an ideal energy extraction target.

\section{Reaction Mechanism}

\subsection{Step 1: Neutron Production}

A free proton with kinetic energy $E_k \geq 1.293$ MeV moves through
dense hydrogen gas. Upon approaching a hydrogen atom, the proton is
attracted to the bound electron:

\begin{equation}
  F = -\frac{e^2}{4\pi\epsilon_0 r^2} \quad (\text{attractive, no Coulomb barrier})
\end{equation}

When the proton and electron interact at sufficiently close range,
inverse beta decay proceeds:

\begin{equation}
  p + e^- \rightarrow n + \nu_e \qquad \Delta mc^2 = -1.293\,\text{MeV}
  \label{eq:capture}
\end{equation}

\subsection{Step 2: Chain Reaction}

After Eq.~\eqref{eq:capture}, the original hydrogen nucleus is
released as a new free proton. This proton repeats the process on the
next hydrogen atom, forming a chain:

\begin{align}
  p^{(1)} + \text{H}_1 &\rightarrow n_1 + \nu_e + p^{(2)} \tag{C1} \\
  p^{(2)} + \text{H}_2 &\rightarrow n_2 + \nu_e + p^{(3)} \tag{C2} \\
  &\vdots \notag \\
  p^{(i)} + \text{H}_i &\rightarrow n_i + \nu_e + p^{(i+1)} \tag{Ci}
\end{align}

Each step consumes one hydrogen atom and produces one neutron and
one free proton that sustains the chain.

\subsection{Step 3: Energy Extraction via Lithium-6}

Produced neutrons are directed into a Lithium-6 blanket:

\begin{equation}
  n + {}^6\text{Li} \rightarrow {}^4\text{He} + {}^3\text{H}
  + 4.78\,\text{MeV}
  \label{eq:licapture}
\end{equation}

This reaction is exothermic, emits no gamma radiation, and
the kinetic energy of products thermalizes into usable heat.

\section{Energy Balance}

\begin{table}[h!]
\centering
\caption{Energy budget per complete reaction cycle}
\vspace{0.5em}
\begin{tabular}{lcc}
\toprule
\textbf{Process} & \textbf{Energy} & \textbf{Direction} \\
\midrule
$p + e^- \rightarrow n + \nu_e$                               & 1.293 MeV & Input  \\
$n + {}^6\text{Li} \rightarrow {}^4\text{He} + {}^3\text{H}$ & 4.780 MeV & Output \\
\midrule
\textbf{Net per cycle}                                         & \textbf{3.487 MeV} & \textbf{Output} \\
\bottomrule
\end{tabular}
\end{table}

\begin{equation}
  \boxed{E_\text{net} = 4.780 - 1.293 = 3.487\,\text{MeV per neutron}}
\end{equation}

\subsection{Chain Sustainability}

The chain sustains itself when each liberated proton retains
sufficient energy to trigger the next capture. Ionization losses
over one mean free path $\lambda$ must not reduce $E_k$ below
1.293 MeV:

\begin{equation}
  E_k^{(i+1)} = E_k^{(i)} - 1.293\,\text{MeV} - E_\text{loss}(\lambda)
  \geq 1.293\,\text{MeV}
\end{equation}

The mean free path at pressure $P$ and temperature $T$ is:

\begin{equation}
  \lambda = \frac{k_B T}{\sqrt{2}\,\pi d_H^2\,P}
\end{equation}

where $d_H \approx 120$ pm. At $P = 100$ atm, $\lambda \approx 2$ nm,
minimizing ionization losses and favoring chain propagation.

\section{Comparison with Conventional Approaches}

\begin{table}[h!]
\centering
\caption{Comparison with related nuclear energy mechanisms}
\vspace{0.5em}
\begin{tabular}{lccc}
\toprule
\textbf{Feature} & \textbf{D--T Fusion} & \textbf{Fission} & \textbf{This Work} \\
\midrule
Coulomb barrier    & High      & None (n)  & \textbf{None} \\
Fuel               & D, T      & U-235     & H$_2$, Li-6   \\
Chain mechanism    & No        & Yes       & Yes           \\
Radioactive waste  & Low       & High      & Low           \\
Energy per event   & 17.6 MeV  & 200 MeV   & 3.487 MeV     \\
\bottomrule
\end{tabular}
\end{table}

\section{Open Challenges}

\begin{enumerate}

  \item \textbf{Energy threshold.}
    The initiating proton must carry $E_k \geq 1.293$ MeV. A
    particle accelerator or high-energy ion source is required to
    start the chain.

  \item \textbf{Ionization losses.}
    The free proton loses energy traversing hydrogen gas. High
    pressure operation ($P > 50$ atm) is necessary to keep $\lambda$
    short enough to sustain the chain.

  \item \textbf{Neutrino energy loss.}
    Each event emits a neutrino carrying $\sim$0.511 MeV that
    escapes unrecovered. This is included in the energy balance.

  \item \textbf{Neutron capture efficiency.}
    Neutrons must reach the Li-6 blanket before escaping or
    thermalizing elsewhere. Neutron reflectors and optimized
    geometry are required.

\end{enumerate}

\section{Proposed Experimental Test}

\begin{enumerate}

  \item Compress hydrogen gas to $P > 50$ atm in a chamber
    surrounded by a Li-6 neutron capture blanket.

  \item Inject a proton beam at $E_k = 5$--20 MeV into the gas.

  \item Measure neutron flux at the Li-6 blanket versus hydrogen
    pressure. A superlinear dependence on pressure would indicate
    chain propagation.

  \item Measure heat output via calorimetry. Net output exceeding
    beam energy input would confirm the energy-positive cycle.

\end{enumerate}

\section{Conclusion}

We have proposed a Coulomb-barrier-free neutron production mechanism
based on proton--electron capture in dense hydrogen gas, with energy
extraction via Lithium-6 neutron capture. The key features are:

\begin{itemize}
  \item No Coulomb barrier (proton--electron attraction).
  \item Self-propagating chain reaction via liberated protons.
  \item Net energy output: 3.487 MeV per reaction cycle.
  \item Abundant, low-cost fuel: hydrogen and Lithium-6.
\end{itemize}

Whether the chain can sustain itself against ionization losses is
the central open question, directly testable with existing
accelerator technology.

\section*{Acknowledgments}

The author thanks the nuclear and LENR physics communities whose
published work informed this theoretical development.

\begin{thebibliography}{9}

\bibitem{widom2006}
  A. Widom and L. Larsen,
  ``Ultra low momentum neutron catalyzed nuclear reactions on metallic
  hydride surfaces,''
  \textit{Eur. Phys. J. C}, vol. 46, pp. 107--111, 2006.

\bibitem{fleischmann1989}
  M. Fleischmann and S. Pons,
  ``Electrochemically induced nuclear fusion of deuterium,''
  \textit{J. Electroanal. Chem.}, vol. 261, pp. 301--308, 1989.

\bibitem{pdg2022}
  R. L. Workman et al. (Particle Data Group),
  ``Review of Particle Physics,''
  \textit{Prog. Theor. Exp. Phys.}, 083C01, 2022.

\bibitem{knoll2000}
  G. F. Knoll,
  \textit{Radiation Detection and Measurement}, 3rd ed.
  Wiley, 2000.

\end{thebibliography}

\end{document}
